\gostunnumberedsection{Заключение}{Заключение}
\label{sec:conclusion}
В первой главе выполнен систематический обзор методов решения задач классификации при дисбалансе классов. Проанализированы data-level и algorithm-level подходы, метрики качества, а также специфика прикладных постановок с табличными данными: экономическая асимметрия ошибок, distribution drift и ограничения на deployment моделей. Сопоставление существующих решений показало, что ни один универсальный метод не доминирует на всех наборах данных; для промышленных задач критичны согласованность метрик с бизнес-KPI и reproducibility протокола сравнения. Сформулированы требования к предметно-ориентированному подходу, включающие явную модель стоимости ошибок, адаптивный rebalancing и процедуру threshold calibration.

Во второй главе разработана концепция и формальное описание предметно-ориентированного подхода, интегрирующего этапы feature engineering и data preprocessing, model training и decision rule по threshold с учётом бизнес-ограничений. Приведено теоретическое обоснование границ применимости и описаны программная реализация и архитектура experimental pipeline. Определён план экспериментов: набор baseline-методов, метрики (PR-AUC, recall at fixed precision, экономический gain) и правила статистического сравнения. Детальное наполнение главы результатами проектирования и параметрами алгоритма будет представлено по мере завершения исследовательской части работы.

В третьей главе (по мере проведения вычислительных экспериментов) будут представлены описание наборов данных, hyperparameter tuning моделей, численные результаты и сравнение с базовыми и современными методами. В качестве внешней проверки воспроизводимости предполагается использовать открытые наборы TabReD с обязательным temporal test \cite{rubachev2025tabred}; по согласованию допускается включение отдельных задач GraphLand для оценки методов на graph data при сильном дисбалансе и temporal shift \cite{bazhenov2025graphland}. Основной промышленный контекст задаётся проектом партнёра (кредитная организация): дефолт и связанные редкие события в табличных и sequence modeling постановках \cite{smirnov2021habrcard,smirnov2022habrbki,meacca2022habraccount,nickimpark2023habrunified}. Планируется анализ robustness к distribution shift, оценка reproducibility (фиксация версий данных, кода и конфигураций в репозитории) и error analysis с точки зрения предметной области. На текущем этапе подготовки работы экспериментальные результаты являются ожидаемыми и будут уточнены после завершения расчётов на полном объёме данных партнёра.

Общая оценка достижения цели. Цель работы - разработка и экспериментальная проверка предметно-ориентированных подходов к задачам с дисбалансом классов - достигается поэтапно: теоретическая и проектная части структурированы, программный каркас и протокол экспериментов определены; окончательная оценка достижения цели будет дана после получения и анализа полного массива экспериментальных данных.

Научная новизна и практическая значимость результатов связаны с системной формализацией предметных ограничений в едином конвейере обучения и принятия решений, ориентированном на продуктовые KPI. Практическая значимость подтверждается возможностью внедрения в процессы модельной разработки при соблюдении требований информационной безопасности и регламентов валидации моделей.

Рекомендации по внедрению и использованию включают: (1) согласование целевых метрик с владельцами продукта до начала training; (2) использование единого experiment repository с версионированием данных и конфигураций; (3) мониторинг не только predictive performance, но и операционных KPI (false positive rate, нагрузка на операторов); (4) периодическая recalibration threshold при distribution drift.

Направления дальнейших исследований: расширение подхода на multi-class постановки и learning with rejection; интеграция с online learning и bandit-style threshold optimization; исследование fairness constraints при cost-sensitive обучении; применение tabular foundation models при дисбалансе классов.

Рекомендации по завершению работы (ожидаемый итог диссертации): (1) явная формулировка предметной постановки (стоимость ошибок, ограничения на false positive rate, требования к interpretability и latency scoring); (2) воспроизводимый evaluation protocol сравнения baseline и предлагаемого метода на обезличенных данных партнёра и на открытом temporal split; (3) отдельный блок про robustness к temporal shift и мониторинг performance degradation (аналогично MLOps-практикам); (4) сжатый набор выводов: в каких условиях предметно-ориентированная связка ``metrics - rebalancing - threshold calibration'' даёт выигрыш над универсальными схемами.
